% Folder 119, batch 16 — archivists' pages 301–313.
% Pass: Opus 5.5 (claude-opus-5-5), 2026-10-03 — first pass, unchecked against the pages by a human.
%
% Contents of the batch, item by item (special rules for folder 119, issue #45):
%  - p. 301: last page (« - 14 - ») of a typed list of lectures 1990–1991,
%    the end of Serre's Collège de France résumé 1990–91. Third-party, printed:
%    \page plus a note only.
%  - pp. 302, 307, 309: blank archivists' divider sheets. No \page.
%  - pp. 303–306: « Curriculum vitae de Alexandre Grothendieck », typescript,
%    4 pp. (« -2- » … « -4- »), undated (internal evidence: written after
%    September 1970 and before the 1972 CNRS letter of p. 308). His text,
%    with one word in ink (« du », p. 303). TRANSCRIBED in full, including the
%    typed slip pasted at the head of p. 303 (struck « et » in ink) and the
%    « Principales publications » [1]–[19].
%  - p. 308: CNRS form letter (form « AV- 20.I.71 »), handwritten date
%    « le 4/7/72 », signed S. Gigon, returning his « Titres et Travaux ».
%    Third-party: \page plus a note only.
%  - pp. 310–313: second copy of the same CV. Not re-transcribed: \page plus
%    notes recording its variants only — the manuscript correction on p. 311
%    and the manuscript entry [20] on p. 313, both transcribed with \add.
% The batch map supplied with the task agrees with the pages.
\documentclass[11pt,a4paper]{article}
\input{../preamble/grothendieck}

\folder{119}
\batch{16}
\pages{301}{313}
\foldertitle{Esquisse d’un programme [dossier constitué en vue d'une candidature au CNRS (1984)] : copies de tapuscrits et de tapuscrits annotés (1972-1974, 1978-1979, 1984, 1990-1991, s.d.), tirés à part (1971), notes et copies de note manuscrites (1986, s.d.), lettres (1968, 1972, 1991).}
\dating{1968-1991}
\watermark{Édition de démonstration}

\begin{document}

\page{301}
\note{Dernière page (« - 14 - ») d'un document dactylographié de J.-P. Serre : fin de son résumé de cours et travaux au Collège de France pour 1990--1991 (liste d'exposés, novembre 1990 -- juin 1991). Texte d'un tiers, déjà imprimé : non transcrit.}

\section{CURRICULUM VITAE DE ALEXANDRE GROTHENDIECK}

\page{303}
\note{En tête de la page est collée une bande dactylographiée, d'une autre machine que le curriculum vitae ; le « et » y est biffé à l'encre.}

Vous trouverez ci-joint l'exposé des titres \struck{et} travaux des candidats à une direction de recherche

Né le 28 mars 1928 à Berlin, de mère allemande et de père apatride, émigré de Russie en 1921, mes parents émigrent d'Allemagne en 1933, participent à la révolution espagnole; je les rejoins en mai 1939. Mes parents sont internés, d'abord mon père en 1939, puis ma mère en 1940 avec moi. Mon père est déporté du camp de Vernet en août 1942 pour Auschwitz et est resté disparu; ma mère meurt en 1957 des suites d'une tuberculose contractée au camp de concentration. Je reste près de deux ans dans des camps de concentration français, puis suis recueilli par une maison d'enfants du "Secours suisse" au Chambon-sur-Lignon, où je termine mes études de lycée en 1945. Etudes de licence (mathématiques) à Montpellier 1945-48, auditeur libre à l'Ecole Normale Supérieure à Paris en 1948-49, où je suis le premier séminaire Cartan sur la théorie des faisceaux, et un cours de Leray \add{du} Collège de France sur la théorie de Schauder du degré topologique dans les espaces localement convexes. De 1949 à 1953 je poursuis des recherches à Nancy sur les espaces vectoriels topologiques, comme élève de J. Dieudonné et de L. Schwartz, aboutissant à ma thèse de doctorat en 1953, sur la théorie des produits tensoriels topologiques et des espaces nucléaires, publiée dans les "Memoirs of the American Mathematical Society". Je passe alors deux ans à l'Université de Saõ Paulo (Brésil), où je continue et mène à leur aboutissement naturel certaines recherches liées aux produits tensoriels topologiques [6,7], mais en même temps, sous l'influence de J. P. Serre, commence à me familiariser avec des questions de topologie algébrique et d'algèbre homologique. Ces dernières continueront à m'occuper jusqu'à aujourd'hui, et sont encore très loin d'être menées à leur terme. Ce sont elles qui m'occuperont surtout pendant l'année 1955 passée à l'Université du Kansas (USA); j'y développe une théorie commune pour la théorie de Cartan-Eilenberg des foncteurs dérivés des foncteurs de modules et la théorie de Leray-Cartan de la cohomologie des faisceaux [8], et développe des notions de "cohomologie non commutative" dans le contexte des faisceaux et des espaces fibrés à

\page{304}
faisceau structural, qui trouveront leur cadre naturel quelques années plus tard avec la théorie des topos (aboutissement naturel du point de vue faisceautique en topologie générale) [16, SGA 4].

A partir de 1956 je suis resté en France, à l'exception de séjours de quelques semaines ou mois dans des universités étrangères. De 1950 à 1958 j'ai été chercheur au CNRS, avec le grade de directeur de recherches en 1958. De 1959 à 1970 j'ai été professeur à l'Institut des Hautes Etudes Scientifiques. Ayant découvert à la fin de 1969 que l'IHES était subventionné depuis trois ans par le Ministère des Armées, et après des essais infructueux pour inciter mes collègues à une action commune sans équivoque contre la présence de telles subventions, je quitte l'IHES en septembre 1970.

Depuis 1959 je suis marié à une française, et je suis père de quatre enfants. Je suis apatride depuis 1940, et ai déposé une demande de naturalisation française au printemps 1970.

Depuis 1956 jusqu'à une date récente, mon intérêt principal s'est porté sur la géométrie algébrique. Mon intérêt pour la topologie, la géométrie analytique, l'algèbre homologique ou le langage catégorique a été constamment subordonné aux multiples besoins d'un vaste programme de construction de la géométrie algébrique, dont une première vision d'ensemble remonte à 1958. Ce programme est poursuivi systématiquement dans [16, 17], d'abord dans un isolement relatif, mais progressivement avec l'assistance d'un nombre croissant de chercheurs de valeur. Il est loin d'être achevé à l'heure actuelle. L'extraordinaire crise écologique que nous aurons à affronter dans les décades qui viennent rend peu probable qu'il le sera jamais. Elle nous imposera d'ailleurs une perspective et des critères de valeur entièrement nouveaux, qui réduiront à l'insignifiance ("irrelevance") beaucoup des plus brillants progrès scientifiques de notre siècle, dans la mesure où ceux-ci restent étrangers au grand impératif évolutionniste de notre temps : celui de la survie. Cette optique s'est imposée à moi avec une force croissante au cours de discussions avec de nombreux collègues sur la responsabilité sociale des scientifiques, occasionnées par ma situation à l'IHES depuis la fin de 1969. Elle m'a conduit en juillet 1970 à m'associer à la

\page{305}
fondation d'un mouvement international et interprofessionnel "Survivre", et à consacrer aux questions liées à la survie une part importante de mon énergie. Dans cette optique, la seule valeur de mon apport comme mathématicien est de me permettre aujourd'hui, grâce à l'estime professionnelle et personnelle acquise parmi mes collègues, de donner plus de force à mon témoignage et à mon action en faveur d'une stricte subordination de toutes nos activités, y compris nos activités de scientifiques, aux impératifs de la survie, et à la promotion d'un ordre stable et humain sur notre planète, sans lequel la survie de notre espèce ne serait ni possible, ni désirable.

A. GROTHENDIECK.

\subsection{PRINCIPALES PUBLICATIONS}

\emph{Espaces Vectoriels Topologiques}

[1] Critères de compacité dans les espaces fonctionnels généraux, Amer. J. 74 (1952), p. 168-186.

[2] Sur certains espaces de fonctions holomorphes, J. Crelle 192 (1953), p. 35-64 et 77-95.

[3] Espaces Vectoriels Topologiques, Notes polyc., Saõ Paulo (1954), 240 p.

[4] Sur les espaces (F) et (DF), Summa Bras. 3 (1954), p. 57-123.

[5] Produits tensoriels topologiques et espaces nucléaires, Mem. AMS, n° 16 (1955), 329 p.

[6] Résumé de la théorie métrique des produits tensoriels topologiques, Bull. Saõ Paulo 8 (1953), p. 1-79.

[7] La théorie de Fredholm, Bull. SMF, 84 (1956), p. 319-384.

\emph{Topologie et algèbre homologique}

[8] Sur quelques points d'algèbre homologique, Tohoku M.J., 9 (1957), p. 119-221.

[9] Théorèmes de finitude pour la cohomologie des fasceaux\note{sic}, Bull. SMF, 84 (1956), p. 1-7.

\emph{Géométrie analytique}

[10] Sur la classification des fibrés holomorphes sur la sphère de Riemann, Amer. J., 79 (1957), p. 121-138.

[11] Techniques de construction en géométrie analytique, Sem. H. Cartan, 13 (1960/61), exposés 7 à 17.

\page{306}
\emph{Géométrie algébrique}

[12] La théorie des classes de Chern, Bull. SMF 86 (1958), p. 137-154.

[13] Sur une note de Mattuck-Tate, J. Crelle 200 (1958), p. 137-154.
\note{la pagination de [13] répète celle de [12] ; telle quelle sur la page.}

[14] The cohomology theory of abstract algebraic varieties, Proc. Int. Congress, Edinburgh (1958), p. 103-118.

[15] Elements de Géométrie Algébrique (rédigés avec la coll. de Jean DIEUDONNE), Chap. I-IV, publ. Math. IHES (1960/67), env. 1 800 pages.

[16] Séminaires de Géométrie Algébrique (SGA 1, …, 7), IHES, 1960/69, env. 4 000 pages (en cours de réédition chez Springer, Lecture Notes) :
\begin{itemize}
  \item SGA 1 Théorie du Groupe Fondamental
  \item SGA 2 Cohomologie locale et Théorèmes de Lefschetz locaux et globaux
  \item SGA 3 (en coll. avec M. Demazure) Schémas en Groupes des Topos et Cohomologie étale des Schémas
  \item SGA 5 Cohomologie l-adique et fonctions L
  \item SGA 6 (en coll. avec J. Berthelot et J.L. Illusie) Théorie des Intersections et Théorèmes de Riemann-Roch
  \item SGA 7 Groupe de Monodromie en Géométrie Algébrique
\end{itemize}
\note{la ligne « SGA 4 » manque : « Théorie des Topos et Cohomologie étale des Schémas », titre de SGA 4, est dactylographié comme suite de la ligne SGA 3, dont l'article initial « Théorie » est tombé ; telle quelle sur la page.}

[17] Un théorème sur les homomorphismes de schémas abéliens, Invent. Math. 2 (1966), p. 59-78.

[18] Dix exposés sur la cohomologie des schémas (en coll. avec J. Giraud, S. Kleiman, M. Raynaud, J. Tate), North Holland, 1968.

[19] Catégories cofibrées additives et complexe cotangent relatif, Lecture Notes in Maths., Springer n° 79 (1968), 167 pages.

\page{308}
\note{Lettre-type du CNRS (Division B, Bureau 2 B, Chercheurs), datée à la main « le 4/7/72 », signée S. Gigon pour le Directeur général, adressée à Grothendieck à Massy : à l'issue de la session du Comité national, elle lui retourne des exemplaires de sa notice de titres et travaux. Texte d'un tiers : non transcrit.}

\page{310}
\note{Seconde copie du curriculum vitae (pp. 310--313), identique au texte des pp. 303--306 et non retranscrite ; seules ses variantes sont relevées. Page 310 : même bande collée, même « du » à l'encre ; aucune variante.}

\page{311}
\note{Variante : correction manuscrite, en écriture fine, dans l'interligne et la marge droite de la phrase « L'extraordinaire crise écologique … rend peu probable qu'il le sera jamais. » L'ajout commence au-dessus de « écologique que » et se poursuit sous la ligne ; « nous » est biffé, « aurons à affronter » barré d'un trait fin, une virgule est ajoutée après « viennent ». « que » n'est pas biffé sur la page. Un signe d'encre sous le point après « actuelle » reste d'interprétation incertaine.}

… Il est loin d'être achevé à l'heure actuelle. L'extraordinaire crise écologique \add{de civilisation dans laquelle nous sommes engagés à l'heure actuelle, et qui \uncertain{devra} se résoudre} que \struck{nous aurons à affronter} dans les décades qui viennent\add{,} rend peu probable qu'il le sera jamais. …

\page{312}
\note{Seconde copie, page « -3- » : aucune variante.}

\page{313}
\note{Seconde copie, page « -4- » : la liste est prolongée à la main, sous [19], par une vingtième référence.}

\add{[20] Platitude d'une adhérence schématique et lemme de Hironaka généralisé (avec H. Seydi), manuscripta math., 5, 323-339, (1971)}

\end{document}
